\subsection{Native Coroutine Bridge}

\texttt{async def} creates a coroutine object at call time (body not executed). \texttt{await} suspends the coroutine frame until an awaitable completes. An event loop schedules coroutine progress cooperatively.

\subsubsection{\texttt{async def} Functions as Native Coroutine Object Factories}

\begin{verbatim}
async def fetch():
    return 42

coro = fetch()
print(f"type: {type(coro)}")
print(f"cr_frame before await: {coro.cr_frame}")
\end{verbatim}

Possible output:

\begin{verbatim}
type: <class 'coroutine'>
cr_frame before await: <frame at 0x..., file "...", line 1, code fetch>
\end{verbatim}

Like generators, coroutines hold a \texttt{cr\_frame} on the heap until completion, then release it.

\subsubsection{\texttt{await} as Structured Suspension Over Awaitable Objects}

\texttt{await expr} evaluates \texttt{expr}, obtains an awaitable, and suspends the current coroutine frame until the awaitable signals completion. Local state in the frame persists across suspension points.

\subsubsection{Event Loops as External Schedulers for Coroutine Progress}

The event loop (\texttt{asyncio.run}, \texttt{loop.run\_until\_complete}) drives coroutine execution: resume one task until its next \texttt{await}, switch to another ready task, repeat. Scheduling is cooperative and single-threaded by default; parallelism requires multiple loops/processes or explicit executor offloading.

Native coroutines replace generator-based coroutine patterns (\texttt{@asyncio.coroutine}, \texttt{yield from}) with dedicated \texttt{coroutine} objects and \texttt{await} syntax, but the underlying frame-suspension model remains the same: heap-allocated \texttt{PyFrameObject} state preserved across voluntary yields.
